\section{Introduction}
\label{sec:intro}

End-to-end driving policies are compared by benchmark scores. A score summarises what the policy's
planned trajectories achieved on logged scenes: no collision, staying in the drivable area, making
progress~\cite{dauner2024navsim,cao2025pseudosim}. An investor reads a fund's return figure the same
way, and no investor allocates capital on the return figure alone. They also ask how large a position
the fund takes, whether it manages risk when risk appears, whether it becomes more careful or more
reckless as the market heats up, and how much it is moved by noise. A deployment decision about a
driving policy needs the same four answers. It usually has to reach them from little data, because a
closed-loop simulator for the target domain or a fleet test is exactly what the decision is meant to
save.

We propose a \emph{check-up} that answers these four questions from a few hundred logged frames. Like
a clinical examination it consists of a short list of exams (\cref{tab:exams}). Each exam has a known
normal range and a control, and each result is read as a diagnosis rather than a ranking. Every exam is
a counterfactual edit of real frames. We remove a pedestrian from the ego corridor, or we re-render the
frame at night, and read the change in the policy's planned trajectory. An independent detector
verifies every edit. The central exam is a \emph{hazard-sensitivity curve}. It measures how much safer
the whole planned trajectory is because the policy can see the pedestrian (in separation, contact and
time to proximity), against a model-independent hazard level. A specificity term applies a
minimal-change principle: when no reaction is needed, the plan should not move.

\paragraph{Diagnosis}
Applied to six end-to-end policies, the check-up separates quantities that a benchmark score merges
(\cref{tab:report}). The policies take very different position sizes: their plans cover
\NumExpLo--\NumExpHi{} times the distance the human drove. Yet none of them manages the pedestrian
risk. Seeing the pedestrian makes the plan safer by \NumHSlo--\NumHShi{} on a $[-1,1]$ scale, and in
five of six policies this benefit does not grow with the hazard. Collisions with the pedestrian's
logged future therefore follow position size and not perception. A night rendering moves every policy
more than the pedestrian does. We look inside the policies to explain the diagnosis, although this
paper does not attempt a repair. The pedestrian leaves a trace in the features but is not decoded by the policies' own semantic
heads; the planner reads a grid too coarse to locate it; and its plan ends before the pedestrian.

\paragraph{How much data}
A check-up is only useful if it is cheap. We fit how each estimate's spread shrinks with the number of
frames. Exposure, over-reaction and the lighting verdict follow the $1/\sqrt{n}$ law closely and are
settled within about 50 frames. Hazard sensitivity and the direction of the lighting effect need
hundreds of frames. We then pre-registered seven predictions from a diagnosis made on 88 Boston frames
and tested them on Singapore. Orderings and verdicts transfer; point values do not
(\cref{tab:prereg}).

\paragraph{Contributions}
\begin{enumerate}\itemsep2pt
\item A small-sample check-up for driving policies. It consists of verified counterfactual edits, a
      hazard-sensitivity curve built on whole-trajectory separation, time to proximity and a
      minimal-change principle, and a reading of the results as a risk profile.
\item A diagnosis of six end-to-end policies, with an explanation from inside the networks of why none
      of them manages pedestrian risk.
\item An account of what small data can and cannot establish: sample-size laws for every exam and a
      pre-registered cross-city transfer test.
\end{enumerate}
